\section{Experiments and Evaluation Setup}
\label{sec:exp}

\subsection{Datasets and Clinical Inter-Patient Split}
\label{sec:data_split}
The empirical validation relies on two internationally recognized PhysioNet benchmark databases:
\begin{enumerate}
    \item \textbf{MIT-BIH Arrhythmia Database (MITDB) \cite{moody2001}:} Contains 48 ambulatory two-lead recordings (30 minutes each, sampled at $f_s = 360$\,Hz) from 47 subjects. The modified limb lead II (MLII) is extracted. Following AAMI standards, four recordings containing paced rhythms (102, 104, 107, 217) are excluded, leaving 44 clinically valid records.
    \item \textbf{MIT-BIH Noise Stress Test Database (NSTDB) \cite{moody1984}:} Provides clean baseline wander (BW), muscle artifact (MA), and electrode motion (EM) noise recordings acquired at 360\,Hz from healthy volunteers.
\end{enumerate}

\paragraph*{Clinical Inter-Patient Protocol (de Chazal DS1/DS2)}
To prevent patient-specific data leakage, we enforce the gold-standard inter-patient division established by de Chazal \emph{et al.} \cite{dechazal2004}:
\begin{itemize}
    \item \textbf{Dataset 1 (DS1 - Development Set, 22 records):} Used exclusively for model parameter optimization. To guarantee strict reproducibility, records \texttt{116}, \texttt{118}, and \texttt{223} are fixed as the validation set ($7{,}279$ beats), while the remaining 19 records form the training set ($43{,}299$ beats).
    \item \textbf{Dataset 2 (DS2 - Independent Test Set, 22 records):} Comprises \textbf{49{,}298 independent beats} from 22 entirely unseen patients, serving as the blind diagnostic test set.
\end{itemize}
Table~\ref{tab:ds_split} details the exact beat distributions mapped to the AAMI EC57 diagnostic classes: Normal and Bundle Branch Block beats ($\mathbf{N}$), Supraventricular Ectopic beats ($\mathbf{S}$), and Ventricular Ectopic beats ($\mathbf{V}$).

\begin{table}[!htbp]
\caption{Beat Distributions Across the Inter-Patient DS1/DS2 Split (AAMI Classes N, S, V)}
\label{tab:ds_split}
\centering
\resizebox{\columnwidth}{!}{%
\begin{tabular}{@{}llrrrr@{}}
\toprule
Partition & Records Included & N & S & V & Total \\
\midrule
DS1 Training & 19 records & 39{,}335 & 774 & 3{,}190 & 43{,}299 \\
DS1 Validation & 3 records (116, 118, 223) & 6{,}511 & 170 & 598 & 7{,}279 \\
\textbf{DS2 Test Set} & \textbf{22 unseen records} & \textbf{44{,}241} & \textbf{1{,}837} & \textbf{3{,}220} & \textbf{49{,}298} \\
\bottomrule
\end{tabular}%
}
\end{table}

\subsection{Experimental Configurations and Baselines}
\HSL{} is systematically evaluated against five prominent state-of-the-art ECG denoising baselines implemented under the identical data preparation and noise protocol:
\begin{enumerate}
    \item \textbf{DW-CNN (Jin \emph{et al.}, 2024 \cite{jin2024}):} Deep wavelet convolutional neural network replacing spatial pooling with DWT operations ($333.8$\,K parameters).
    \item \textbf{DNN-DAN (Xiong \emph{et al.}, 2016 \cite{xiong2016}):} Deep neural network denoising autoencoder ($230.7$\,K parameters).
    \item \textbf{FCN (Chiang \emph{et al.}, 2019 \cite{chiang2019}):} Fully convolutional denoising autoencoder ($738.9$\,K parameters).
    \item \textbf{LiWave (Nguyen \emph{et al.}, 2026 \cite{nguyen2026liwave}):} Lightweight 1D wavelet-residual U-Net ($82.1$\,K parameters).
    \item \textbf{DeepFilter (Romero \emph{et al.}, BSPC 2021 \cite{romero2021}):} Multi-branch dilated convolutional filter for baseline wander removal ($68.7$\,K parameters).
\end{enumerate}

\subsection{Training Hyperparameters}
Models were implemented in PyTorch 2.4. In Stage~1, \HSL{} was trained with Adam ($\text{lr} = 10^{-3}$, batch size 16, 50 epochs, early stopping patience 10) on an NVIDIA RTX GPU. In Stage~2, \RRN{} was optimized using AdamW ($\text{lr} = 10^{-3}$ for classical layers, $\text{lr} = 2 \times 10^{-4}$ for quantum circuit parameters, weight decay $10^{-4}$, batch size 64, early stopping patience 15). For edge latency benchmarks, single-thread forward inference was evaluated on a standard Intel Core CPU at batch size 1.

\subsection{Evaluation Metrics}
\paragraph*{Waveform Reconstruction Fidelity}
With $x_0[n]$ representing the clean reference, $x[n]$ the contaminated input, and $\hat{y}[n]$ the denoised signal over segment length $L$:
\begin{align}
\mathrm{RMSE} &= \sqrt{\frac{1}{L} \sum_{n=0}^{L-1} \big(x_0[n] - \hat{y}[n]\big)^2}, \\
\mathrm{PRD}  &= 100 \times \frac{\sqrt{\sum_{n=0}^{L-1} (x_0[n] - \hat{y}[n])^2}}{\sqrt{\sum_{n=0}^{L-1} x_0^2[n]}}\ \ (\%), \\
\Delta\mathrm{SNR} &= 10\log_{10}\frac{\sum x_0^2}{\sum (x_0 - \hat{y})^2} - 10\log_{10}\frac{\sum x_0^2}{\sum (x_0 - x)^2}\ \ (\mathrm{dB}).
\end{align}

\paragraph*{Clinical Morphological Fiducials (NeuroKit2 \cite{makowski2021})}
Using the validated physiological library NeuroKit2, we quantify clinical landmark errors: R-peak amplitude error $e_R = |\hat{A}_R - A_R|$ (mV), QRS duration deviation $\Delta\mathrm{QRS}_{\mathrm{dur}}$ (ms), and QT interval error $\Delta\mathrm{QT}_{\mathrm{int}}$ (ms).

\paragraph*{Classification Metrics}
For each diagnostic class $c \in \{N, S, V\}$, we measure Sensitivity ($\mathrm{Se} = \frac{\mathrm{TP}}{\mathrm{TP} + \mathrm{FN}}$), Positive Predictivity ($+\mathrm{P} = \frac{\mathrm{TP}}{\mathrm{TP} + \mathrm{FP}}$), and $F_1\text{-score} = \frac{2 \cdot \mathrm{Se} \cdot {+\mathrm{P}}}{\mathrm{Se} + {+\mathrm{P}}}$, alongside overall Accuracy ($\mathrm{Acc}$) and Macro-averaged F1 ($\mathrm{Macro\text{-}F1}$).
